\section{Evidence classes and reconstruction without source glyphs}
\label{sec:supp_l3}
Registered on 2026-10-23, after the material generator was closed and before any run; amendments are dated below. \textbf{Question.} With the target string known, can the text be written correctly, in place and stably when no reliable glyph exists, from the carrier alone? \textbf{Evidence classes.} The frozen localiser of \cref{sec:method} was run once on the intact benchmark, restricted to matches inside the hint box: a clip is \emph{glyphs} when the target is matched on at least 20\% of its frames, \emph{pseudo-text} when the detector finds text boxes inside the hint on at least 20\% of the frames but nothing matches, and \emph{glyph-free} otherwise. Of the 526 clips, \lthreeGlyphs{} are glyphs, \lthreeLtwo{} pseudo-text and \lthreeLthree{} glyph-free. In our selected benchmark, every clip contains an OCR-detectable text-like region, although \lthreeLtwo{} clips lack a reliable transcription; this is a property of this generator, this selection rule and this detector, and a missing detection would not have meant that the text had visually vanished. In the run of \cref{tab:main}, \method reads \lthreeGlyphsBoneRead\% of the glyphs clips and \lthreeLtwoBoneRead\% of the pseudo-text clips (through its detection-box fallback), \steer \lthreeGlyphsSteerRead\% and \lthreeLtwoSteerRead\%. The registered natural glyph-free experiment and its routing regression are therefore empty; both were cancelled before any output existed.

\textbf{The branch.} Inputs are the clip, the frame-0 hint box (the same box the fallback uses; never the per-frame boxes) and the target. The carrier is tracked by SIFT homographies around the hint (the estimator of the \emph{track} flag; a clip is \emph{locatable} when it succeeds on at least 80\% of the frames), its extent is scanned outward from the hint on the middle frame with the rules of the carrier-fit module (colour, texture and image-edge boundaries) using the hint's core as the carrier colour, and the text box is laid out deterministically inside that extent (height at most the hint's and 0.6 of the usable height, the string's own aspect ratio, at most 0.86 of the usable width, centred on the hint). Ink is a contrast colour against the carrier, stroke 0.12 of the text height, no glow; the hint region is erased once on the middle frame and carried by the carrier homographies; the glyph canvas of \cref{sec:method} is rendered through the same homographies. The branch never sees glyph masks, OCR boxes, letter heights or source colours.

\textbf{Experiment A, synthetic glyph-free clips.} For the \lthreeGlyphs{} glyphs clips the tracked source quads are the geometry truth; the input is the clip with that quad erased (LaMa once on the middle frame, carried by the quad homographies: \method's own erase), the frame-0 hint and the target. Comparator: the frozen code with its static-hint fallback enabled (\texttt{GL\_FALLBACK=mask}); the deployed configuration ships an unlocated clip unchanged and is reported too (it located \lthreeADepN{} of the 455 clips and read \lthreeADepRead\% of those). Criteria, all required: readability at most one point below the fallback; per-frame IoU with the truth higher (paired CI above zero) and centre drift lower (CI below zero); BG PSNR $\geq 41$\,dB; flicker within $0.05$ and jitter within $0.2$ of the fallback; judged on the locatable clips. \textbf{Result} (\cref{tab:l3}): on the \lthreeALoc{} locatable clips the branch reads \lthreeACarRead\% against \lthreeABfbRead\% (\lthreeAWins{} vs.\ \lthreeALosses{} discordant, $p=\lthreeAP$), keeps the background (\lthreeACarBG\,dB) and lands closer to the old text's centre (\lthreeACarOffset{} vs.\ \lthreeABfbOffset{} text heights, drift \lthreeACarDrift{} vs.\ \lthreeABfbDrift), but its box is not closer to the truth (IoU \lthreeACarIoU{} vs.\ \lthreeABfbIoU, paired \lthreeAIoUMean{} [\lthreeAIoULo, \lthreeAIoUHi]; it renders the text at \lthreeACarHeightRatio{} of the original height) and it is less stable (flicker \lthreeAFlickerMean{} [\lthreeAFlickerLo, \lthreeAFlickerHi], jitter \lthreeAJitterMean{} [\lthreeAJitterLo, \lthreeAJitterHi] above the fallback): \lthreeAPassed{} of \lthreeACriteria{} criteria, \textbf{not passed}. The truth is the old text's footprint, not the optimal layout of the new string, so the IoU criterion is a consistency check only; we report it as registered. The \lthreeANotLoc{} non-locatable clips are reported separately. \textbf{Post hoc} (second run, registered after the result). The true text motion is small (287 of 455 clips move by less than 0.3\,px per frame) and the tracker's velocity error (0.22\,px per frame) equals the motion itself; frames without a SIFT estimate fell back to the static hint, producing jumps. Interpolating across failed frames, smoothing over 9 frames and holding the reference quad when the carrier hardly moves gives readability \lthreeAtwoRead\%, flicker \lthreeAtwoFlicker{} (within the limit) and jitter \lthreeAtwoJitter{} (\lthreeAtwoJitterPMean{} [\lthreeAtwoJitterPLo, \lthreeAtwoJitterPHi] above the fallback, still outside), IoU \lthreeAtwoIoU: the temporal instability was partly attributable to noisy carrier tracking under near-static camera motion. The registered verdict stands.

\textbf{Leakage audit of the synthetic clips.} On the middle frame, the old ink (the colour segmentation of \cref{sec:metrics} on the intact frame) is compared with the other pixels of the old quad after erasure. By the rule fixed before the audit (gradient ratio above 1.5 or a mean luminance difference above 5 levels) \lthreeAuditFlagOrig{} of \lthreeAuditN{} clips are flagged; the gradient ratio has median \lthreeAuditGradMedian{} (no glyph-shaped edges) and the luminance difference median \lthreeAuditDlumMedian{} levels. A revised rule comparing the old ink with a ring around it and with the ink--luminance correlation flags \lthreeAuditFlagRev{} (\lthreeAuditBoth{} in common); the four clips with the strongest correlation show cleanly erased text and a carrier structure (window frames, box edges) inside the old quad; the reader's similarity to the target on the erased inputs is 0 in median. Both counts are kept. Splitting by the original rule, the branch reads \lthreeAuditSuspectCar\% against \lthreeAuditSuspectBfb\% on the flagged clips and \lthreeAuditLowCar\% against \lthreeAuditLowBfb\% on the others: its advantage does not depend on residual cues, but the erased carriers keep low-frequency brightness differences in about half of the clips, which we record as a limitation of the synthetic protocol.

\textbf{Amendment B$'$, natural pseudo-text clips} (registered once the classes were known, before any output of the branch). The branch against \method's deployed outputs on the \lthreeLtwo{} pseudo-text clips, with the criteria of the cancelled natural experiment: readability not significantly worse, no more broken clips, BG PSNR $\geq 41$\,dB, flicker and jitter not worse in CI. Readability \lthreeBCarRead\% vs.\ \lthreeBBoneRead\% (\lthreeBWins{} vs.\ \lthreeBLosses{} discordant, $p=\lthreeBP$), broken \lthreeBCarBroken{} vs.\ \lthreeBBoneBroken, BG \lthreeBCarBG\,dB, flicker and jitter within CI: \textbf{passed}; the branch repairs all \lthreeBFailures{} clips \method fails and loses one. The readability gain is not statistically significant at this size and is bought with appearance: the distance rises from \lthreeBBoneMad{} to \lthreeBCarMad{} (\lthreeBMadMean{} [\lthreeBMadLo, \lthreeBMadHi]), since the branch renders flat text without \method's material modules. A diagnostic run that renders \method's own per-frame layout through the branch's erase and renderer reads \lthreeDiagOcrRead\%, so the gain comes from the layout source (carrier extent vs.\ detection boxes: \lthreeDiagWins{} vs.\ \lthreeDiagLosses, $p=\lthreeDiagP$) and not from the erase; the comparison is nevertheless branch-level.

\textbf{Reproducibility note.} The same frozen code, re-run on the pseudo-text clips in the second environment (OpenCV 4.11, CUDA 12.6, V100; the deployed outputs come from OpenCV 4.12, CUDA 12.8, A40), reads \lthreeBBrerunRead\% instead of \lthreeBBoneRead\% (\lthreeRerunWins{} vs.\ \lthreeRerunLosses{} discordant) with BG PSNR \lthreeBBrerunBG{} instead of \lthreeBBoneBG\,dB. Clip by clip, the decoded frames are identical, the erase masks agree (IoU \xenvMaskIoU), the layout quads agree (\xenvGeomPx\,px) and the OCR quads coincide, but the font matcher chooses a different typeface in \xenvFontDiff{} of \xenvN{} clips and the OCR instance decision differs in a few (letters equal in \xenvLettersSame, gate in \xenvGateSame), so the glyph alphas agree only to IoU \xenvAlphaIoU{} and no output frame is bit-identical (PSNR \xenvPsnrIn\,dB inside the text window, \xenvPsnrOut{} outside). A same-environment repeat is bit-exact. The front end is deterministic within an environment and sensitive to library versions at the font-matching step; B$'$ is therefore reported against the deployed outputs and against the same-environment re-run (\lthreeBrWins{} vs.\ \lthreeBrLosses{} discordant, BG \lthreeBrBGMean{} [\lthreeBrBGLo, \lthreeBrBGHi]).

\textbf{Verdict.} B1 remains the system. The branch is an exploratory capability with a stated boundary: it writes readable text where the deployed system writes nothing, and it does so from the carrier alone, but it has not shown that it can do so without losing material and some temporal stability, and it is not merged. The registered routing regression was not triggered; the exploratory routing study follows.

\textbf{Pseudo-text routing impact study (exploratory; not the registered regression).} All 526 clips in one evaluation run: \method (B1) against a tree that keeps B1 on the \lthreeGlyphs{} glyphs clips and routes the \lthreeLtwo{} pseudo-text clips to the branch (the 455 shared outputs are the same files, with no discordant reading). Readability \cexBoneRead\% $\to$ \cexCexRead\% (\cexWins{} vs.\ \cexLosses{} discordant clips, $p=\cexP$), net repairs \cexBoneNet{} $\to$ \cexCexNet, broken clips \cexBoneBroken{} $\to$ \cexCexBroken; BG PSNR \cexBoneBG{} $\to$ \cexCexBG\,dB (\cexBGMean{} [\cexBGLo, \cexBGHi]); flicker and jitter unchanged (\cexFlickerMean{} [\cexFlickerLo, \cexFlickerHi], \cexJitterMean{} [\cexJitterLo, \cexJitterHi]); appearance distance \cexBoneMad{} $\to$ \cexCexMad{} over all clips (\cexMadMean{} [\cexMadLo, \cexMadHi]) and \cexBoneLtwoMad{} $\to$ \cexCexLtwoMad{} on the routed clips, rising in every material (most on screens, metal and wood). With the post-hoc tracking on the routed clips the temporal metrics improve slightly (flicker \cexCexAFlicker, jitter \cexCexAJitter) and nothing else changes. The routing buys about one point of readability and two fewer broken clips with a measurable loss of appearance on the routed clips and a small loss of background fidelity; by the judgement fixed before this run, B1 remains the system and the branch is not merged.
\IfFileExists{tables/tab_l3.tex}{\begin{table}[t]
\centering\small\setlength{\tabcolsep}{4pt}
\caption{Carrier-grounded reconstruction without source glyphs (supplement). A: the 445 synthetic glyph-free clips whose carrier the branch can track (one evaluation run; the post-hoc row is a second run). B$'$: the 71 pseudo-text clips of the benchmark (one run). Readability by the first reader; BG PSNR in dB; temporal metrics as in Tab.~1.}
\label{tab:l3}
\resizebox{\linewidth}{!}{%
\begin{tabular}{lcccc}
\toprule
Row & Read\,$\uparrow$ & BG\,$\uparrow$ & Flicker\,$\downarrow$ & Jitter\,$\downarrow$\\
\midrule
\multicolumn{5}{l}{\textit{A: synthetic glyph-free clips, locatable carriers (445)}}\\
Erased input & 0.0 & -- & 1.81 & 7.16\\
Frozen \method, static-hint fallback & 81.6 & 41.8 & 2.26 & 8.81\\
Carrier branch (registered) & 92.1 & 41.6 & 2.73 & 12.36\\
Carrier branch, post-hoc tracking (second run) & 92.1 & 41.6 & 2.30 & 10.41\\
\midrule
\multicolumn{5}{l}{\textit{B$'$: natural pseudo-text clips (71)}}\\
Original & 32.4 & -- & 2.05 & 9.38\\
\method (deployed outputs) & 90.1 & 42.4 & 1.75 & 7.86\\
\method re-run, same code & 93.0 & 41.0 & 1.76 & 7.91\\
Carrier branch & 98.6 & 41.2 & 1.67 & 7.87\\
\bottomrule
\end{tabular}}
\end{table}
}{}

